\documentclass[instructions.tex]{subfiles}
\begin{document}
 
\chapter{De la vérité des peines de l’enfer.}
 
Il est aussi vrai qu’il y a un enfer, qu’il est vrai qu’il y a un Dieu. C’est nier un Dieu, dit saint Augustin, que de ne pas croire qu’il pardonne au pécheur pénitent. C’est aussi nier la Divinité que de ne pas croire un Dieu vengeur du crime.
 
Dire qu’il n’y a ni paradis ni enfer, c’est dire que Dieu ne se soucie point de l’homme, qu’il le conduit au hasard, qu’il l’abandonne à lui-même, qu’il voit ses désordres et qu’il les approuve ; c’est dire que Dieu est sans sagesse, sans sainteté et sans justice. Penser de la sorte, c’est avoir de fausses idées de la Divinité, ce n’est pas connoître Dieu, c’est être athée.\footnote{Recordare, Jesu pie, quòd sum causa tuæ viæ; ne me perdas illà die. Prose.}
 
Les incrédules et les impies ont beau raisonner sur ce point si sérieux ; toujours doivent-ils croire, 1. qu’il peut y avoir un enfer : car oseroient-ils disputer ce pouvoir à la toute-puissance du Très-Haut ? 2. Ils doivent de plus croire que l’enfer est au moins douteux : car quelle assurance ont-ils du contraire ? prétendent-ils avoir plus de lumières que tout l’univers, et que les plus savans génies qui l’ont cru ? 3. Ils doivent enfin croire que l’enfer est certain : la foi le dit, et la raison le démontre.
 
Oui, l’enfer est un article de la foi divine : Jésus-Christ en son Évangile, dans un seul chapitre, répète jusqu’à six et sept fois que le feu de l’enfer ne s’éteindra jamais\footnote{Marc 9.}. Un feu, dit le Seigneur, a été allumé dans ma fureur ; il brûlera jusqu’au fond des enfers ; j’y rassemblerai tous les maux, j’y épuiserai tous les traits de ma colère\footnote{Deut. 32.}.
 
Le prophète Isaïe, ce grand homme, plus grand encore par la force et l’élévation de son génie, que par sa naissance royale, cet homme plein de l’esprit de Dieu, parle ainsi : Le Seigneur a dit : Écoutez, vous qui vous éloignez de moi, connoissez ma puissance : les pécheurs seront dans l’épouvante ; le tremblement saisira les hypocrites. Qui de vous pourra habiter dans un feu dévorant, dans des ardeurs éternelles ?
 
Mais quand la foi et l’autorité de Dieu ne nous obligeraient pas à le croire, la raison nous le démontre. Les païens, par la lumière naturelle, l’ont connu, et les plus éclairés n’en ont jamais douté. Nous connoissons, par notre raison, qu’il y a un Dieu ; cette même raison nous dit qu’il est saint, qu’il est juste ; il doit donc défendre et punir le vice : étant Dieu, il doit le punir en Dieu ; ses châtimens doivent donc être bien au-dessus de tout ce que les hommes peuvent imaginer. Voilà ce qu’une raison éclairée apprend à tout homme de bon sens. Faites l’esprit résolu tant qu’il vous plaira, il sera toujours vrai qu’il y a un enfer que vous devez le croire, et qu’on le croira toujours.
 
\secundo Si l’on ne croit pas cette vérité, c’est parce qu’on est vicieux. Otez l’orgueil d’un esprit pointilleux, et l’impureté d’un cœur dépravé, et bientôt il sera convaincu qu’il y a un enfer. Les méchans voudraient qu’il n’y en eût point, et que leurs crimes fussent impunis ; faut-il s’étonner qu’ils s’étourdissent pour se persuader qu’il n’est point ? S’ils disent qu’il n’y a point d’enfer, ils sont d’impudens menteurs, leur propre conscience les dément. Vous osez le dire, mais n’osez pas le penser, dit saint Augustin ; votre conscience vous fait sentir le contraire : vous avez beau en étouffer les remords, elle vous fait toujours savoir qu’on vous en menace, et vous fait sentir que vous le méritez.
 
Je ne pense qu’à tirer parti de la vie, disent-ils, sans me mettre en peine de ce qui arrivera. O raisonnement insensé ! Une éternité de tourmens, est-ce donc une chose indifférente. Quand l’enfer serait aussi douteux qu’il est certain, la pensée de ce malheur devrait déjà vous effrayer. Un mal douteux est toujours à craindre. Prendriez-vous une liqueur, si vous doutiez qu’elle fût empoisonnée, ou un chemin, si vous doutiez d’y être assassiné, votre doute ne fût-il fondé que sur le témoignage d’un imbécile ? Ainsi, l’enfer ne fût-il que douteux, ce serait déjà le comble de l’imprudence de ne pas travailler à l’éviter.
 
C’est parce que vous craignez de faire pénitence, qu’il vous plaît de douter de l’enfer. Voilà, dit saint Augustin, encore un trait de folie, de craindre une pénitence de quelques jours, et de ne pas craindre une éternité de feux ; si vous prenez tant de précautions contre la mort, contre une maladie qui ne dure que peu de temps, pourquoi ne vous précautionnez-vous pas pour éviter une mort et des supplices qui dureront toujours\footnote{Timent cruciatum temporalem, et non timent pœnas ignis æterni ; timent modicum mori, et non timent æternùm mori. Serm. 18. de Verb.} ? Vous ne vous souciez pas de l’enfer, pourvu que vous jouissiez de vos plaisirs ici-bas. Or, n’est-ce pas avoir perdu l’esprit, que d’acheter si cher, non pas des plaisirs, mais des repentirs éternels ? A la bonne heure, disent certains impies : si je suis un jour en enfer, je m’en consolerai : du moins je n’y serai pas seul. Quoi, malheureux ! vous vous feriez une peine de demeurer en retraite pendant un mois avec de saintes personnes ; et vous regardez comme une consolation d’habiter, non pas un mois, mais pendant l’éternité, avec des hommes réprouvés de Dieu, plongés dans un étang de feu, qui souffrent des tourmens horribles ; qui, dans la rage, dans un désespoir furieux, se maudissent et se déchirent ! Penser de la sorte, ce n’est pas force d’esprit, mais fascination, frénésie, fureur.
 
Concluons que celui qui ne croit et ne craint pas l’enfer, a perdu la foi et la raison : la foi, puisque l’autorité de Dieu oblige de le croire ; la raison, puisqu’elle nous dicte qu’un Dieu souverainement juste doit défendre le vice et le punir en Dieu. Voilà tout ce que l’univers croit, ce que tous les hommes pieux et les plus grands génies ont toujours cru et croiront toujours. Malheur à ceux qui plaisantent sur des vérités si terribles ! Laissez railler, laissez plaisanter les impies, les voluptueux et les fous, qui rient aujourd’hui de ce qui les fera pleurer un jour amèrement ; laissez parler ceux qui veulent se perdre. Pour vous, devenez sage à leurs dépens, et sauvez-vous.
 
\section{Résolutions}
 
\primo Je me dirai de temps en temps : Je crois l’enfer sur la parole de Dieu ; je l’ai souvent mérité, puisqu’un seul péché mortel suffit pour s’en rendre digne : jusqu’à quand attendrai-je donc de me convertir, et de commencer une vie vraiment pénitente ? 

\secundo Dès ce moment je renonce à tout ce qui m’empêche de travailler à éviter le souverain malheur. 

\prayer{Soyez béni, ô mon Dieu, de la patience que vous avez exercée envers moi ; ah ! faites que j’en profite sans plus tarder, pour mettre ordre à ma conscience, et assurer mon salut.}
 
\end{document}
